\chapter{Stand der Technik} \label{chpt:Stand_der_Technik_Main}

Ein RISC-V Komitee hat im Jahre 2015 mit der ersten Version der RISC-V Vector Extension (\gls{RVV}) begonnen. Im Jahre 2021 wurde das Release 1.0 ratifiziert und ist jetzt in einen finalen Zustand (Frozen) überführt worden. RISC-V wird durch die (\gls{RVV}) um Register für Vektoren und explizite Anweisungen zur Manipulation dieser Vektoren erweitert.

%Diese führte dazu, dass bereits Versuche unternommen worden sind, kleine lokale Large-Language-Modelle auf RISC-V CPUs auszuführen. Dabei zeigt sich, dass sich kleine LLMs tatsächlich mit ? Token pro Sekunde unterhalten können.

In der RISC-V Vector Extension Spezifikation \cite{rvv_specification} wird das Konzept des Strip-Mining beschrieben und hervorgehoben. Die Software lädt den Cache mit einem Teil einer oder mehrerer Vektoren voll und beginnt den Cache dann abzuarbeiten, in dem der Teil der Vektoren dann wiederum teilweise in den CPU geladen, verarbeitet und danach wieder aus dem CPU zurück in den Speicher übertragen wird. Das Strip-Mining wiederholt diese Abläufe, bis die Aufgabe erledigt ist, also alle Element der Vektoren bearbeitet sind. Das Beispiel in Kapitel 6.4. namens ``Example of stripmining and changes to SEW'' aus \cite{rvv_specification} verdeutlicht dieses Vorgehen.

%Diese zwei Abläufe (Cache-Frame laden und Register laden) werden dann solange wiederholt, bis der gesamte Vektor abgearbeitet ist. Ein Hinweis auf dieses Strip-Mining Vorgehen findet sich in der RISC-V Vector Extension Spezifikation. zum Strip-Mining.

Die Alternative zum Hardware unterstützen Strip-Mining wäre es, jedes einzelne Vektorelement einzeln in ein Register der CPU zu laden und dann mit einem Element des zweiten Vektors zu verrechnen. Dabei wird für jedes Element ein Laden und Speicher aus dem Speicher notwendig und man bezahlt mit einem hohen Verlust an Geschwindigkeit für die im Gegenzug triviale Implementierung.

Der Vorteil der \gls{RVV} Register ist es, dass diese Vektorregister mehrere Elemente parallel auf eine Instruktion hin zusammenrechnen können. Der CPU verfährt dann in einem echten \gls{SIMD}-Betriebsmodus (Single Instruction Multiple Data) und spart ein vielfaches an Taktzyklen im Vergleich zur sequentiellen Ausführung.

%Es sind auch komplett neuartige Designs denkbar. Wenn der Cache-Speicher über mehrere Kanäle an den CPU angebunden ist, könnte der CPU die Vektorregister auch parallel über alle Kanäle befüllen, mehrer Register parallel verrechnen und die Ergebnisse parallel wieder in den Cache zurückliefern.

% https://github.com/flame/blis/blob/master/docs/KernelsHowTo.md

Das Prinzip des Strip-Mining wird auch in mathematischen Bibliotheken verwendet. Eine bekannte Bibliothek ist blis \cite{BLIS1}. blis verwendet den Begriff Kernel oder Microkernel. Ein Kernel ist Code, der auf einem Cache-Frame ausgeführt wird, also den Code darstellt, den die CPU sehr schnell ausführen kann. Optimierte Anweisung für verschiedene CPU Architekturen werden hauptsächlich in dem Code verwendet, der dem Kernel zuzurechnen ist um Zeit zu sparen.

Um die Berechnungen im Kernel herum gibt es weiteren Code, der über Schleifen nacheinander alle Daten in den Kernel lädt, bis die Berechnung vollständig abgeschlossen ist.

Es gibt sehr starke Parallelen zwischen dem Strip-Mining und dem Kernel-Ansatz der blis Bibliothek. Sollte blis für RISC-V portiert werden, ist es sinnvoll, dass RISC-V das Prinzip des Strip-Mining ermöglicht. Dies wurde durch die RISC-V Vector Extension für Vektoren bereits erreicht.

%blis dokumentiert die Matrixmultiplikation unter dem Stichwort GEMM. GEMM steht für General Matrix Multiply. Der gemm kernel wird für die GEMM Algorithmen verwendet.

Für die RISC-V Matrixberechnung gibt es zum jetzigen Zeitpunkt keine ratifizierte Extension. Innerhalb der RISC-V Arbeitsgruppen werden unterschiedliche Ansätze diskutiert. Die Ansätze werden im Kapitel zu den theoretischen Grundlagen beschrieben.

Interessanterweise gibt es eine chinesische Firma namens Stream Computing die sich zur Mission gesetzt hat, high-end RISC-V Chips zum Einsatz in Rechenzentren zu entwickeln. Bereits im Jahre 2022 konnte Stream Computing ein eigenes Arbeitsergebnis vorweisen \footnote{https://riscv.org/blog/stream-computing-risc-v-matrix-extension-open-source-project-upgrades-to-version-0-5-supporting-vectormatrix-implementation/}. Dieses Arbeitsergebnis wurde von Stream Computing selbst ``RISC-V Matrix ISA Specification'' in der Version v0.1 genannt. Die Version v0.5 wurde im Oktober 2024 erstellt. Es handelt sich aber trotz des Namens nicht um die offizielle Version, die das RISC-V Gremium veröffentlichen wird! Stream Computing hat ebenfalls die notwendigen Anpassungen zum LLVM Compiler-Framework beigetragen!

Wie tragbar die Ergebnisse sind lässt sich schwer einschätzen ohne eine eingehende Evaluierung des Quellcodes durchgeführt zu haben. Es ist jedoch interessant zu sehen, welche Innovationskraft und Arbeitsgeschwindigkeit die chinesischen Industrie hat. Während sich das RISC-V Gremium wieder auf eine mehrjährige Reise begibt um eine Matrix-Extension zu erzeugen, wie es bereits bei der Vektor-Extension erfolgt ist, hat eine chinesische Firma vermeintlich bereits ein Design und eine Toolchain zu diesem Design erstellt.

% TODO Verweis: https://riscv.org/blog/stream-computing-risc-v-matrix-extension-open-source-project-upgrades-to-version-0-5-supporting-vectormatrix-implementation/

Innerhalb des NEORV32 CPU Umfeldes gibt es eine Arbeit von Qizal Arsalan \cite{QArsalan_NEORV32_MATRIX_CFU} bei der der NEORV32 CPU um Matrixberechnungen erweitert wird. Diese Bachelorarbeit unterscheidet sich allerdings hinsichtlich der Herangehensweise stark von Qizal Arsalans Arbeit. Der NEORV32 CPU besitzt im Design vorgesehene Erweiterungspunkte namens CFS und CFU. Das CFS (Custom Functions Subsystem) erlaubt es Hardware-Beschleuniger in den NEORV32 CPU einzubringen, die über Speicheradressen (Memory Mapping) aus einer Softwareanwendung mit Daten versorgt werden können. CFU (Custom Function Unit) erlaubt es einzelne Instruktionen zu ergänzen um den CPU um kleinschrittige Operationen zu erweitern.

Qizal Arsalan untersucht den Ansatz der Matrixmultiplikation mit CFS und CFU. In dieser Bachelorarbeit werden Matrix-Instruktionen in den NEORV32 Chip eingebracht, die nicht über das CFS oder CFU Feature eingebunden werden sondern über die Execution-Engine, also der State-Machine, des NEORV32 Chips selbst. Der Grund ist, dass die Speicheraddressabhängigkeit von CFS und CFU nicht konform mit einer Matrix-Extension-Spezifikation des RISC-V Gremiums sein können und nicht sein werden.

Bei CFS und CFU handelt es sich um Systeme, die nicht Teil der RISC-V Spezifikation sind. Diese Bachelorarbeit versucht die Ergebnisse einer Matrix-Extension im RISC-V Umfeld zu antizipieren und kann daher CFS und CFU nicht verwenden.